(1) 令\(x=E-E_0\)与\(a=\Gamma/2\)，有\(S=e^{2i\delta_{\rm bg}}(x-ia)/(x+ia)\)，由于分子分母互为复共轭，\(|S|=1\)。取连续相位分支写\(S=e^{2i\delta(E)}\)，共振相位的导数为\(d\delta_{\rm res}/dE=a/(x^2+a^2)\)；在跨越极点的解析延拓中，相位连续增加\(\pi\)，不能因为主值的反正切跳跃而把导数符号弄错。
(2) 单通道Wigner–Smith延迟\(\tau_W=-i\hbar S^*(dS/dE)=2\hbar\,d\delta/dE\)。背景相移若近似常量，\(\tau_W(E)=\hbar\Gamma/[(E-E_0)^2+(\Gamma/2)^2]\)，在\(E_0\)处峰值\(4\hbar/\Gamma\)，半高全宽\(\Gamma\)。该延迟定义的是散射相位对能量的导数，峰值不应不加区分地等同于概率驻留寿命\(\hbar/\Gamma\)。
(3) 若共振能向左、右两端口衰减，分宽\(\Gamma_L,\Gamma_R>0\)，在无其他漏失通道时\(\Gamma=\Gamma_L+\Gamma_R\)。标准共振透射率\(T(E)=\Gamma_L\Gamma_R/[(E-E_0)^2+(\Gamma/2)^2]\)，峰值\(T(E_0)=4\Gamma_L\Gamma_R/(\Gamma_L+\Gamma_R)^2\le1\)，等号仅在\(\Gamma_L=\Gamma_R\)；不对称耦合使峰值下降却不改变能量守恒。
(4) 当\(\Gamma\to0^+\)，共振能量宽度变窄、峰值延迟按\(1/\Gamma\)变长，说明长寿命准束缚态与窄共振之间的联系。若背景相移随能量缓变，应在\(\tau_W\)中增加\(2\hbar\,d\delta_{\rm bg}/dE\)，它可正可负；孤立Lorentz表达仅适用于背景可近似常量且无重叠共振的窗口。
